\chapter{Categorías fibradas y descenso}
\label{VI}
\marginpar{145}%

\setcounter{section}{-1}
\section{Introducción}

Contrariamente a lo anunciado en la introducción de la exposición
anterior, resultó imposible efectuar el descenso en la categoría de los
preesquemas, aun en casos particulares, sin haber desarrollado antes con
suficiente cuidado el lenguaje del descenso en las categorías generales.

La noción de «descenso» proporciona el marco general para todos los
procedimientos de «pegado» de objetos y, por consiguiente, de «pegado» de
categorías. El caso más clásico de pegado corresponde al dato de un espacio
topológico~$X$ y de un recubrimiento de~$X$ por abiertos~$X_i$; si para cada~$i$
se da un espacio fibrado (digamos)~$E_i$ sobre~$X_i$, y para cada par~$(i,j)$
un isomorfismo~$f_{ji}$ de~$E_i|X_{ij}$ sobre~$E_j|X_{ij}$ (donde se pone
$X_{ij}=X_i\cap X_j$), de modo que estos isomorfismos satisfacen una conocida
condición de transitividad
(escrita abreviadamente~$f_{kj}f_{ji}=f_{ki}$), se sabe que existe un espacio
fibrado~$E$ sobre~$X$, determinado salvo isomorfismo por la condición de que
existan isomorfismos~$f_i\colon E|X_i\isomto E_i$ que satisfagan las relaciones
$f_{ji}=f_jf_i^{-1}$ (con el abuso de notación habitual). Sea~$X'$ la suma
topológica de los~$X_i$, que es, por tanto, un espacio fibrado sobre~$X$ (es decir,
dotado de una aplicación continua~$X'\to X$). El dato de los~$E_i$ puede
interpretarse de manera más concisa como un espacio fibrado~$E'$ sobre~$X'$,
y el dato de los~$f_{ji}$ como un isomorfismo entre las dos imágenes inversas
(por las dos proyecciones canónicas)~$E''_1$ y~$E''_2$ de~$E'$ sobre
$X''=X'\times_X X'$; la condición de pegado puede escribirse entonces como
una identidad entre isomorfismos de los espacios fibrados~$E'''_1$ y~$E'''_3$
sobre el producto fibrado triple~$X'''=X'\times_X X'\times_X X'$ (donde
$E'''_i$ designa la imagen inversa de~$E'$ sobre~$X'''$ por la proyección
canónica de índice~$i$). La construcción de~$E$ a partir de~$E'$ y de~$f$ es
un caso típico
\marginpar{146}%
de procedimiento de «descenso». Por otra parte, partiendo de un espacio
fibrado~$E$ sobre~$X$, se dice que
\ifthenelse{\boolean{orig}}{$X$}{$E$} es «localmente trivial», de fibra~$F$,
si existe un recubrimiento abierto~$(X_i)$ de~$X$ tal que los~$E|X_i$ sean
isomorfos a~$F\times X_i$, o, lo que es equivalente, tal que la imagen inversa
$E'$ de~$E$ sobre~$X'=\coprod_i X_i$ sea isomorfa a~$X'\times F$.

Así, tanto la noción de «pegado» de objetos como la de «localización» de una
propiedad están ligadas al estudio de ciertos tipos de «cambios de base»
$X'\to X$. En geometría algebraica, muchos otros tipos de cambio de base, y
en particular los morfismos~$X'\to X$ fielmente planos, deben considerarse
como correspondientes a un procedimiento de «localización» relativo a los
preesquemas, u otros objetos, «sobre»~$X$. Este tipo de localización se utiliza
tanto como la simple localización topológica (que, por otra parte, es un caso
particular de aquel). Lo mismo sucede (en menor medida) en geometría analítica.

La mayoría de las demostraciones, que se reducen a verificaciones, se omiten
o solo se esbozan; cuando procede, precisamos los diagramas menos evidentes
que aparecen en una demostración.
